% 张三慧《大学物理学》第4版，第二卷第13章《电势》
% 教材 PDF 物理页 61--64；原印刷页44--47。原图13.16--13.19仍待矢量重绘。
% 原题扫描核对工作稿，未经独立学术评审及版权发布审查。
\section*{第13章\quad 电势}
\subsection*{章末习题（13.1--13.34）}
\paragraph{习题 13.1}
两个同心球面，半径分别为 \(10\,\mathrm{cm}\) 和 \(30\,\mathrm{cm}\)，
小球均匀带有正电荷 \(1.0\times10^{-8}\,\mathrm C\)，
大球均匀带有正电荷 \(1.5\times10^{-8}\,\mathrm C\)。
求离球心分别为（1）\(20\,\mathrm{cm}\)，（2）\(50\,\mathrm{cm}\)
的各点电势（以无穷远为零）。

\paragraph{习题 13.2}
两均匀带电球壳同心放置，半径分别为
\(R_1,R_2\ (R_1<R_2)\)，已知内外球壳之间的电势差为 \(U_{12}\)，
求两球壳间的电场分布。

\paragraph{习题 13.3}
两个同心的均匀带电球面，半径分别为
\(R_1=5.0\,\mathrm{cm},R_2=20.0\,\mathrm{cm}\)，
已知内球面电势 \(\varphi_1=60\,\mathrm V\)，
外球面电势 \(\varphi_2=-30\,\mathrm V\)。
（1）求内、外球面上所带电量；
（2）在两个球面之间何处的电势为零？

\paragraph{习题 13.4}
两个同心的球面，半径分别为 \(R_1,R_2\ (R_1<R_2)\)，
分别带有总电量 \(q_1,q_2\)，设电荷均匀分布在球面上，
求两球面的电势及二者之间的电势差。
不管 \(q_1\) 大小如何，只要是正电荷，内球电势总高于外球；
只要是负电荷，内球电势总低于外球。试说明原因。

\paragraph{习题 13.5}
一细直杆沿 \(x\) 轴由 \(x=-a\) 延伸到 \(x=a\)，
杆上均匀带电，线电荷密度为 \(\lambda\)。
计算 \(x\) 轴上 \(x>a\) 各点的电势（零点取在无穷远）。

\paragraph{习题 13.6}
一均匀带电细杆，长 \(l=15.0\,\mathrm{cm}\)，
线电荷密度 \(\lambda=2.0\times10^{-7}\,\mathrm{C/m}\)，求：
（1）细杆延长线上与杆的一端相距 \(a=5.0\,\mathrm{cm}\) 处的电势；
（2）细杆中垂线上与细杆相距 \(b=5.0\,\mathrm{cm}\) 处的电势。
电势零点取在无穷远。

\paragraph{习题 13.7}
求习题 12.18 中两同轴圆筒之间的电势差。

\paragraph{习题 13.8}
一计数管中有一直径为 \(2.0\,\mathrm{cm}\) 的金属长圆筒，
在圆筒轴线上装有一根直径为
\(1.27\times10^{-5}\,\mathrm m\) 的细金属丝。
设金属丝与圆筒间的电势差为 \(1.0\times10^3\,\mathrm V\)，求：
（1）金属丝表面的场强大小；
（2）圆筒内表面的场强大小。

\paragraph{习题 13.9}
一无限长均匀带电圆柱，体电荷密度为 \(\rho\)，
截面半径为 \(a\)。
（1）用高斯定律求柱内外的场强分布；
（2）求柱内外的电势分布，以轴线为电势零点；
（3）画出 \(E-r\) 和 \(\varphi-r\) 的函数曲线。

\paragraph{习题 13.10}
半径为 \(R\) 的圆盘均匀带电，面电荷密度为 \(\sigma_0\)。
求此圆盘轴线上的电势分布：
（1）用例13.4的结果用电势叠加法；
（2）利用例12.6的结果用场强积分法。

\paragraph{习题 13.11}
一均匀带电的圆盘，半径为 \(R\)，面电荷密度为
\(\sigma\)，今将其中心半径为 \(R/2\) 的圆片挖去。
试用叠加法求剩余圆环带在其垂直轴线上的电势分布，
并求中心的电势和电场强度各是多大。

\paragraph{习题 13.12}
（1）一个球形雨滴半径为 \(0.40\,\mathrm{mm}\)，
带有电量 \(1.6\,\mathrm{pC}\)，它表面的电势多大？
（2）两个这样的雨滴碰后合成一个较大的球形雨滴，
这时大雨滴表面的电势又是多大？
\textit{假定合并时体积和净电荷守恒。}

\paragraph{习题 13.13}
金原子核可视为均匀带电球体，
总电量为 \(79e\)，半径为 \(7.0\times10^{-15}\,\mathrm m\)。
求金核表面的电势及其中心的电势。

\paragraph{习题 13.14}
如原书图13.16，两个平行放置的均匀带电圆环，
半径都为 \(R\)，电量分别为 \(+q\) 及 \(-q\)，
两中心间的距离为 \(l\)，并有 \(l\ll R\)。
（1）以两环的对称中心 \(O\) 为坐标原点，
求垂直于环平面的 \(x\) 轴上电势分布；
（2）证明：当 \(x\gg R\) 时，
\(\varphi=\dfrac{ql}{4\pi\varepsilon_0x^2}\)。
\textit{图示：}右环带 \(+q\)，左环带 \(-q\)，间距 \(l\)。

\paragraph{习题 13.15}
用电势梯度法求习题13.5中 \(x\) 轴上 \(x>a\) 各点的电场强度。

\paragraph{习题 13.16}
符号相反的两个点电荷 \(q_1,q_2\) 分别位于
\(x=-b\) 和 \(x=+b\) 两点，
试证 \(\varphi=0\) 的等势面为球面，
并求该球半径和球心的位置。如果二者电量相等，
则此等势面又如何？
\textit{需要区分电荷大小绝对值相等与否，以及退化为平面的情形。}

\paragraph{习题 13.17}
两条无限长均匀带电直线的线电荷密度分别为
\(-\lambda\) 和 \(+\lambda\)，并平行于 \(z\) 轴放置，
和 \(x\) 轴分别相交于 \(x=-a\) 和 \(x=+a\)。
证明：
（1）此系统的等势面和 \(xy\) 平面的交线都是圆，
并求这些圆的圆心位置和半径；
（2）电场线都是平行于 \(xy\) 平面的圆，
并求这些圆的圆心位置和半径。
\textit{零势等势面可以是直线的极限情形。}

\paragraph{习题 13.18}
一次闪电的放电电压大约是 \(1.0\times10^9\,\mathrm V\)，
放电中和的电量约是 \(30\,\mathrm C\)。
（1）求一次放电所释放的能量；
（2）一所希望小学每天消耗电能 \(20\,\mathrm{kW\,h}\)，
上述一次放电释放的电能足够该小学用多长时间？

\paragraph{习题 13.19}
电子束焊接机的电子枪如原书图13.17所示。
K为阴极，A为阳极，其上有一小孔。
阴极发射的电子在阴极和阳极电场作用下聚集成一细束，
以极高的速率穿过阳极上的小孔，射到被焊接的金属上。
已知 \(\varphi_A-\varphi_K=2.5\times10^4\,\mathrm V\)，
并设电子从阴极发射时初速率为零，求：
（1）电子到达金属时具有的动能（用电子伏特表示）；
（2）电子射到金属上时的速率。

\paragraph{习题 13.20}
一边长为 \(a\) 的正三角形，其三个顶点上各放置
\(q,-q,-2q\) 的点电荷，求此三角形重心上的电势。
将一电量为 \(+Q\) 的点电荷由无限远移到重心上，
外力要做多少功？

\paragraph{习题 13.21}
如原书图13.18，在三块互相平行的均匀带电大平面上，
面电荷密度为
\(\sigma_1=1.2\times10^{-4}\,\mathrm{C/m^2}\)，
\(\sigma_2=2.0\times10^{-5}\,\mathrm{C/m^2}\)，
\(\sigma_3=1.1\times10^{-4}\,\mathrm{C/m^2}\)。
A点与平面II相距 \(5.0\,\mathrm{cm}\)，
B点与平面II相距 \(7.0\,\mathrm{cm}\)。
（1）计算A、B两点的电势差；
（2）设把电量 \(q_0=-1.0\times10^{-8}\,\mathrm C\)
的点电荷从A点移到B点，外力克服电场力做多少功？
\textit{图示：}按 \(x\) 从左到右依次为平面I、II、III；
A位于I与II之间，B位于II与III之间，三张均带正电。

\paragraph{习题 13.22}
电子直线加速器的电子轨道沿直线排列的
长约金属圆筒制成，如原书图13.19，
单数圆筒和双数圆筒分别连在两根导线上，接在交流电源的两极上。
由于电势差的正负交替改变，
可以使一个电子团（延续几个微秒）
依次越过两筒间隙时总能被电场加速；
圆筒内没有电场，电子做匀速运动。
这要求各圆筒的长度依次适当加长。
（1）证明要使电子团发出和跨越每个筒间隙时都正好被电势差峰值加速，
圆筒长度应依次为 \(L_1n^{1/2}\)，
其中 \(L_1\) 是第一个筒长，\(n\) 为圆筒序数（非相对论情形）；
（2）设交变电势差峰值为 \(U_0\)，频率为 \(\nu\)，求 \(L_1\)；
（3）电子从第 \(n\) 个筒出来时的动能多大？

\paragraph{习题 13.23}
（1）按牛顿力学计算，把一个电子加速到光速需要多大的电势差？
（2）按相对论的正确公式，静质量为 \(m_0\) 的粒子动能为
\[
 E_k=m_0c^2\left[\frac1{\sqrt{1-v^2/c^2}}-1\right].
\]
试由此计算电子越过上一问所求的电势差时所能达到的
速度是光速的百分之几。

\paragraph{习题 13.24}
假设某一瞬时，氦原子的两个电子正在核的两侧，
它们与核的距离都是 \(0.20\times10^{-10}\,\mathrm m\)。
这种配置状态的静电势能是多少？（把电子与原子核看作点电荷。）

\paragraph{习题 13.25}
根据原子核的 \(\alpha\) 粒子模型，
某些原子核由 \(\alpha\) 粒子的有规则几何排列组成。
例如 \({}^{12}\mathrm C\) 原子核由排列成等边三角形的
3个 \(\alpha\) 粒子组成。
设每对粒子之间的距离为
\(3.0\times10^{-15}\,\mathrm m\)，
求这3个 \(\alpha\) 粒子配置的静电势能（电子伏特）。
（将 \(\alpha\) 粒子看成点电荷。）

\paragraph{习题 13.26}
一条无限长的一维晶体由沿直线交替排列的正负离子组成，
这些离子的电荷大小都是 \(e\)，相邻离子间距为 \(a\)。
求证：（1）每一个正离子所在处的电势为
\(-\dfrac{e}{2\pi\varepsilon_0a}\ln2\)；
（2）任何一个离子的静电势能为
\(-\dfrac{e^2}{4\pi\varepsilon_0a}\ln2\)。
（提示：利用 \(\ln(1+x)\) 展开式。）

\paragraph{习题 13.27}
假设电子是半径为 \(R\)、总电量为 \(e\) 且均匀
分布在其外表面上的球体。
如果静电能等于电子的静止能量 \(m_ec^2\)，
那么以电子的 \(e\) 和 \(m_e\) 表示的电子半径
\(R\) 的表达式是什么？\(R\) 数值上等于多少？
（此 \(R\) 是电子的“经典半径”；现代高能实验确定
电子的电量集中分布在不超过 \(10^{-18}\,\mathrm m\) 的线度范围内。）

\paragraph{习题 13.28}
如果把质子当成半径为 \(1.0\times10^{-15}\,\mathrm m\)
的均匀带电球体，它的静电势能是多少？
这势能是质子的相对论静能的百分之几？

\paragraph{习题 13.29}
铀核带电量为 \(92e\)，
可以近似地认为它均匀分布在一个半径为
\(7.4\times10^{-15}\,\mathrm m\) 的球体内，
求铀核的静电势能。
当铀核对称裂变后，产生两个相同的钯核，
各带电 \(46e\)，总体积和原来一样。
设这两个钯核也可以看成球体，
当它们分离得远时，它们的总静电势能又是多少？
这一裂变释放出的静电能是多少？
按每个铀核都这样对称裂变计算，
\(1\,\mathrm{kg}\) 铀裂变后释放的静电能是多少？
（裂变时释放的“核能”基本上就是这静电能。）

\paragraph{习题 13.30}
一个动能为 \(4.0\,\mathrm{MeV}\) 的 \(\alpha\) 粒子射向金原子核，
求二者最接近时的距离。 \(\alpha\) 粒子的电荷为 \(2e\)，
金原子核的电荷为 \(79e\)，
将金原子核视作均匀带电球体，并认为它保持不动。
已知 \(\alpha\) 粒子质量 \(6.68\times10^{-27}\,\mathrm{kg}\)，
金核质量 \(3.29\times10^{-25}\,\mathrm{kg}\)，
求在此距离时二者的万有引力势能多大？

\paragraph{习题 13.31}
\(\tau\) 子带有与电子一样多的负电荷，
质量是 \(3.17\times10^{-27}\,\mathrm{kg}\)。
它可以穿入核物质而只受电力作用。
设一个 \(\tau\) 子原来静止在离铀核很远的地方，
由于铀核的吸引而向铀核运动。
求它越过铀核表面时的速度多大？
到达铀核中心时的速度又多大？
铀核可看成带有 \(92e\) 的均匀带电球体，
半径为 \(7.4\times10^{-15}\,\mathrm m\)。

\paragraph{习题 13.32}
两个电偶极子的电矩分别为 \(p_1,p_2\)，
相隔距离为 \(r\)，方向相同，都沿着二者连线。
试证明二者的相互作用静电势能为
\[
 U=-\frac{p_1p_2}{2\pi\varepsilon_0r^3}.
\]

\paragraph{习题 13.33}
地球表面上空晴天时电场强度约为 \(100\,\mathrm{V/m}\)。
（1）此电场的能量密度多大？
（2）假设地球表面以上 \(10\,\mathrm{km}\) 范围内
电场强度都是这一数值，
此范围内储存的电场能共是多少 \(\mathrm{kW\,h}\)？

\paragraph{习题 13.34}
按照玻尔理论，氢原子中的电子围绕原子核作圆运动，
维持电子运动的力为库仑力。
轨道的大小取决于角动量，
最小轨道角动量为
\(\hbar=1.05\times10^{-34}\,\mathrm{J\,s}\)，
其他依次为 \(2\hbar,3\hbar\) 等。
（1）证明：如果圆轨道角动量
\(n\hbar\ (n=1,2,3,\ldots)\)，
则其半径
\(r=\dfrac{4\pi\varepsilon_0}{m_ee^2}n^2\hbar^2\)；
（2）证明：在这样的轨道中，
电子的轨道能量（动能与势能）为
\(W=-\dfrac{m_ee^4}{2(4\pi\varepsilon_0)^2\hbar^2n^2}\)；
（3）计算 \(n=1\) 时的轨道能量（以电子伏特表示）。
